% Example presentation using the HSG Beamer theme.
% Compile from THIS folder: pdflatex example.tex  (three passes)
\documentclass[aspectratio=169]{beamer}

% Tell LaTeX where the theme + assets live (one directory up).
\makeatletter\def\input@path{{../}}\makeatother
\usetheme{HSG}

\title{Computer Networks}
\subtitle{Network Security}
\author{Prof. Dr. Bruno Rodrigues}
\institute{University of St.\,Gallen\\Bachelor in Computer Science\\Computer Networks (2,906)}
\date{}
\coursefooter{Computer Networks}
\coverimage{hsg-closing-campus.jpg}    % demo cover; replace with your own photo

% Deck outline — drives the left-margin sidebar progress.
\hsgsetsections{Setup, Symmetric, Asymmetric, Summary}

% Closing-slide contact (sits inside the translucent white box).
\hsgclosingcontact
  {Prof. Dr. Bruno Rodrigues}
  {bruno.rodrigues@unisg.ch}

\begin{document}

% =====================================================================
% 1. COVER
% =====================================================================
\begin{frame}[plain]
  \titlepage
\end{frame}

% =====================================================================
% 2. AGENDA
% =====================================================================
\hsgagenda{
  \item Setup and goals
  \item Symmetric encryption
  \item Asymmetric encryption
  \item Summary
}

% =====================================================================
% Section 1: Setup
% =====================================================================
\hsgsection{Setup}

\begin{frame}
\frametitle{Lecture goals}
\framesubtitle{What we will cover today}

\begin{itemize}
  \setlength\itemsep{0.5em}
  \item \hsgemph{Provide} a general overview of encryption mechanisms.
  \item \hsgemph{Compare} symmetric (AES) and asymmetric (RSA) approaches.
  \item \hsgemph{Explain} session keys (Diffie-Hellman) and TLS.
  \item \hsgemph{Discuss} certificates and digital signatures.
\end{itemize}
\end{frame}

% =====================================================================
% Section 2: Symmetric
% =====================================================================
\hsgsection{Symmetric}

\begin{frame}
\frametitle{Two-column comparison}
\framesubtitle{Symmetric vs. asymmetric, side-by-side}

\begin{columns}[T]
  \column{0.5\textwidth}
    \textbf{Symmetric encryption}
    \begin{itemize}
      \item Single shared key
      \item Fast, suitable for bulk data
      \item Key distribution is the hard part
    \end{itemize}
  \column{0.5\textwidth}
    \textbf{Asymmetric encryption}
    \begin{itemize}
      \item Public / private key pair
      \item Solves key distribution
      \item Orders of magnitude slower
    \end{itemize}
\end{columns}

\vspace{1em}
In practice, \hsgemph{TLS} combines both: asymmetric cryptography
establishes a shared session key, which is then used with a symmetric
cipher for the data stream.
\end{frame}

\hsgstatement{The hard part is key distribution, not the cipher.}

% =====================================================================
% Section 3: Asymmetric
% =====================================================================
\hsgsection{Asymmetric}

\hsgquote
  {The cloud is just somebody else's computer.}
  {Anonymous}

\hsgquoteinverted
  {The best way to predict the future is to invent it.}
  {Alan Kay}

% =====================================================================
% Section 4: Summary
% =====================================================================
\hsgsection{Summary}

\begin{frame}
\frametitle{Takeaways}

\begin{itemize}
  \setlength\itemsep{0.55em}
  \item Symmetric cryptography is fast; asymmetric solves key distribution.
  \item \hsgemph{Certificates} bind public keys to identities;
        \hsgemph{TLS} uses them to bootstrap trust.
  \item \hsgemph{Diffie-Hellman} gives forward secrecy even if long-term
        keys leak later.
\end{itemize}
\end{frame}

% =====================================================================
% CLOSING
% =====================================================================
\hsgclosing

\end{document}
